\chapter{Generative Models}
\label{ch:generative}

\begin{tcolorbox}[feynbox={Sky}{BBg},title={Before and After}]
\textbf{Before this chapter you need:} logarithms (\S\ref{sec:logrules}),
integrals (\S\ref{sec:integrals}), norms (\S\ref{sec:norms}), densities and
expectation (\S\ref{sec:distributions}--\S\ref{sec:expectation}), maximum
likelihood (\S\ref{sec:mle}), gradient descent and minimax
(\S\ref{sec:gradientdescent}, \S\ref{sec:minimax}), neural networks
(Chapter~\ref{ch:neuralnets}).\\
\textbf{After it you will understand:} what a generative model is trying to
learn, what entropy and KL divergence measure, why Jensen--Shannon divergence
gives no usable gradient when the two distributions do not overlap, and how
Kantorovich--Rubinstein duality turns Wasserstein distance into something a
neural network can compute --- which is where $\norm{\nabla D}_2 = 1$ finally
comes from.
\end{tcolorbox}

A generative model learns to produce new samples resembling its training data.
This chapter builds the vocabulary for saying how close two distributions are,
shows why the obvious choice fails, and derives the alternative the models in
this thesis actually use.

%% ─────────────────────────────────────────────────────────────────────────────
\section{Distributions over Data}
\label{sec:datadist}

\situation{
Someone shows you ten thousand photographs of faces. You could not write down a
formula for ``face'', yet you would recognise an eleventh instantly, and you could
sketch something face-like yourself. You have absorbed a distribution without ever
representing it explicitly.
}

\problem{
The set of plausible return series is a vanishingly small corner of all possible
number sequences. To generate new ones we must characterise that corner --- and
writing it down directly is hopeless, so it must be learned.
}

\workedarith{
\textbf{How small is the corner?} A window of 128 daily returns is a point in
$\R^{128}$ (\S\ref{sec:vectors}). Discretise crudely into 10 buckets per day. The
space of sequences has
\[
  10^{128}
\]
elements --- more than the number of atoms in the observable universe, by roughly
$10^{48}$.

BOVESPA's 20-year history supplies about $4{,}953$ windows. The training data
covers a fraction of the space equal to
\[
  \frac{4953}{10^{128}} \approx 5 \times 10^{-125} .
\]

\medskip
\textbf{Yet the task is not hopeless}, because the plausible sequences are not
scattered uniformly. They concentrate: near-zero mean, standard deviation around
$1.6\%$, volatility clustered, tails heavy. Those constraints carve out a
low-dimensional surface inside $\R^{128}$, and it is that surface a generative
model learns.

\medskip
\textbf{A concrete constraint.} Requiring $|\mathrm{ACF}(1)| < 0.03$ alone
eliminates the overwhelming majority of sequences in the space --- a random
sequence has no reason to satisfy it. Each stylised fact removes another vast
portion.
}

\formaldef{
A \textbf{generative model} approximates an unknown data distribution
$p_{\text{data}}$ over a space $\mathcal{X}$, either by
\begin{itemize}
  \item providing a density $p_\theta(x)$ that can be evaluated, or
  \item providing a procedure that draws samples $x \sim p_\theta$.
\end{itemize}
These are different capabilities. GANs give the second without the first: they
sample but cannot report a likelihood.

The \textbf{manifold hypothesis} holds that real data concentrates near a
low-dimensional surface embedded in the high-dimensional space.
}

\analogybreak{
``Learning the distribution'' overstates what any of these models do. They learn
to produce samples that pass the tests we apply. A model can match every metric in
this thesis and still differ from $p_{\text{data}}$ in ways no metric probes ---
which is exactly why the evaluation framework is built around a taxonomy of what
each metric can and cannot see, rather than around a single score.
}

\whyhere{
The manifold picture explains why the shuffled-real control is so awkward. It lies
\emph{exactly} on the correct marginal distribution while sitting nowhere near the
correct manifold, because the manifold is defined by ordering constraints that
permutation destroys. Any distance that ignores ordering rates it perfect.
}

%% ─────────────────────────────────────────────────────────────────────────────
\section{Latent Variables and Sampling}
\label{sec:latent}

\situation{
A pipe organ. A handful of stops and one keyboard produce an enormous variety of
sounds. The controls are few; the outputs are many and richly structured, because
the mechanism between them does the work.
}

\problem{
We need to generate points on a complicated surface in $\R^{128}$. Sampling
directly from a complicated surface is hard. Sampling from something simple and
transforming it is easy --- provided the transformation can be learned.
}

\workedarith{
\textbf{The recipe.} Draw $\mathbf{z}$ from a simple distribution, then apply a
learned map $G$:
\[
  \mathbf{z} \sim \N(0, I) \in \R^{100} , \qquad \mathbf{x} = G(\mathbf{z}) \in \R^{128} .
\]

\medskip
\textbf{Transformation reshapes distributions.} Take $z$ uniform on $[0,1]$ and
$x = z^2$. Uniform draws $0.1, 0.5, 0.9$ map to $0.01, 0.25, 0.81$. The inputs are
evenly spread; the outputs bunch near zero. A simple squaring has produced a
non-uniform distribution.

\medskip
\textbf{Check the density.} For $x = z^2$ with $z$ uniform,
$\Prob(x \leq a) = \Prob(z \leq \sqrt{a}) = \sqrt{a}$, so the density is
$\frac{d}{da}\sqrt{a} = \frac{1}{2\sqrt{a}}$ --- unbounded near $0$, far from
uniform. \checkmark A trivial map has produced a non-trivial distribution; a
neural network can produce a very complicated one.

\medskip
\textbf{Dimensions in this project.} QuantGAN and FinGAN draw
$\mathbf{z} \in \R^{100}$; TimeGAN uses a 24-dimensional noise vector matched to
its hidden dimension. In every case the latent space is far smaller than the
$128$-dimensional output, which is the manifold hypothesis built into the
architecture.
}

\formaldef{
A \textbf{latent variable model} specifies a prior $p(\mathbf{z})$ --- usually
$\N(0,I)$ --- and a map $G_\theta : \mathcal{Z} \to \mathcal{X}$. Sampling is:
draw $\mathbf{z} \sim p(\mathbf{z})$, return $G_\theta(\mathbf{z})$.

The induced density on $\mathcal{X}$ is
\[
  p_\theta(\mathbf{x}) = \int p(\mathbf{z})\,\delta\bigl(\mathbf{x} - G_\theta(\mathbf{z})\bigr)\,d\mathbf{z} ,
\]
which is generally intractable. When $G$ is invertible with a computable Jacobian
determinant (\S\ref{sec:determinant}), the change-of-variables formula gives the
density explicitly --- the basis of normalising flows. GANs impose no such
restriction and forgo the density.
}

\analogybreak{
A low-dimensional latent space is a modelling assumption with teeth. If the true
data manifold has higher dimension than $\mathcal{Z}$, the generator cannot cover
it whatever its training --- some genuine outputs are simply unreachable. Too large
a latent space wastes capacity and can make training unstable. Neither failure is
visible in the loss.
}

\whyhere{
The generator is a function from noise to a return series, and sampling is a
single forward pass. This is why generating a synthetic dataset costs almost
nothing once training is done, and why a collapsed generator is so dangerous: if
$G$ maps every $\mathbf{z}$ to nearly the same output, sampling a million series
yields a million near-copies, and only a diversity check will notice.
}

%% ─────────────────────────────────────────────────────────────────────────────
\section{Maximum Likelihood as a Training Objective}
\label{sec:mletraining}

\situation{
Section~\ref{sec:mle} found a coin's bias by choosing the value that made the
observed flips most probable. The same principle should extend to fitting a
generator: choose parameters making the observed data most probable.
}

\problem{
The principle extends; the computation does not. For a latent-variable generator
the likelihood involves an integral over all latent values that could have
produced a given output, and that integral has no closed form.
}

\workedarith{
\textbf{Where it works.} Fit $\N(\mu, \sigma^2)$ to three observations
$2, 4, 6$. The log-likelihood is
\[
  \ell(\mu,\sigma) = -\frac{3}{2}\ln(2\pi\sigma^2) - \frac{1}{2\sigma^2}\sum_i (x_i - \mu)^2 .
\]
Differentiating in $\mu$ and setting to zero gives $\hat\mu = \bar{x} = 4$;
differentiating in $\sigma$ gives
\[
  \hat\sigma^2 = \tfrac13\bigl[(2-4)^2 + (4-4)^2 + (6-4)^2\bigr] = \tfrac{8}{3} = 2.667 .
\]
Closed form, done.

\medskip
\textbf{Where it fails.} For $\mathbf{x} = G_\theta(\mathbf{z})$ the density at an
observed $\mathbf{x}$ requires
\[
  p_\theta(\mathbf{x}) = \int p(\mathbf{z})\,\delta\bigl(\mathbf{x} - G_\theta(\mathbf{z})\bigr)\,d\mathbf{z} ,
\]
an integral over $\R^{100}$. There is no formula, and numerical integration in 100
dimensions is out of reach.

\medskip
\textbf{Worse than intractable.} If the generator's output lies exactly on a
$100$-dimensional surface inside $\R^{128}$, then any real data point off that
surface has density \emph{exactly zero}, so
$\ln p_\theta(\mathbf{x}) = -\infty$. The objective is not merely hard to compute;
it is infinite for almost every real observation. Section~\ref{sec:jsd} shows this
is the same difficulty in another guise.
}

\formaldef{
\textbf{Maximum likelihood} for a generative model maximises
\[
  \frac{1}{N}\sum_{i=1}^{N}\ln p_\theta(x_i) ,
\]
which as $N \to \infty$ approaches $\E_{p_{\text{data}}}[\ln p_\theta(x)]$ ---
equivalent to minimising $D_{\mathrm{KL}}(p_{\text{data}} \parallel p_\theta)$
(\S\ref{sec:kl}).

Tractable for: explicit densities, autoregressive models, normalising flows.
Intractable for: implicit generators such as GANs, which is why they optimise a
divergence estimated adversarially instead.
}

\analogybreak{
Maximum likelihood is not automatically the right target even where it is
computable. Minimising $D_{\mathrm{KL}}(p_{\text{data}} \parallel p_\theta)$
penalises assigning low probability to real data, but barely penalises assigning
high probability to nonsense. Models trained this way tend to over-spread, hedging
across the whole space --- which for return generation means producing implausible
series rather than admitting ignorance.
}

\whyhere{
This is the structural reason the models here are adversarial. A likelihood-based
generator would need a tractable density, which for the sequence structure at hand
would mean strong architectural restrictions. GANs trade the likelihood away for
freedom of architecture --- and the price is that no model in this thesis can report
how probable a given real series is under it, so evaluation must be entirely
sample-based.
}

%% ─────────────────────────────────────────────────────────────────────────────
\section{Entropy}
\label{sec:entropy}

\situation{
Two weather forecasters. One says ``rain'' every day in a climate where it rains
half the time; the other reports genuine daily odds. Over a year, how many yes/no
questions would you need to pin down the actual weather? That number measures how
uncertain the weather is.
}

\problem{
Before distributions can be compared, we need a way to quantify the uncertainty in
a single one --- and it should have the property that a predictable distribution
scores low and an unpredictable one scores high.
}

\workedarith{
\textbf{A fair coin.} Two outcomes, each $1/2$:
\[
  H = -\left[\tfrac12\log_2\tfrac12 + \tfrac12\log_2\tfrac12\right]
    = -\left[\tfrac12(-1) + \tfrac12(-1)\right] = 1 \text{ bit} .
\]
One yes/no question suffices, and one is needed.

\medskip
\textbf{A biased coin, $p = 0.9$.}
\[
  H = -\bigl[0.9\log_2 0.9 + 0.1\log_2 0.1\bigr]
    = -[0.9(-0.152) + 0.1(-3.322)] = 0.137 + 0.332 = 0.469 \text{ bits} .
\]
Less than half. The outcome is usually heads, so less information arrives per
flip.

\medskip
\textbf{A certain coin, $p = 1$.} $H = -1\log_2 1 = 0$. No uncertainty, no
information.

\medskip
\textbf{A fair die.} Six outcomes each $1/6$:
\[
  H = \log_2 6 = 2.585 \text{ bits} .
\]

\medskip
\textbf{The pattern.} Entropy is maximised by the uniform distribution and
minimised by a certain one. For $n$ equally likely outcomes it is exactly
$\log_2 n$.
}

\formaldef{
The \textbf{Shannon entropy} of a discrete distribution is
\[
  H(p) = -\sum_x p(x)\log p(x) ,
\]
in bits with $\log_2$, in nats with $\ln$. For a continuous density the
\textbf{differential entropy} is
\[
  h(p) = -\int p(x)\ln p(x)\,dx .
\]
Entropy is non-negative in the discrete case, maximised by the uniform
distribution over a finite set, and equal to zero only for a point mass.
}

\analogybreak{
Differential entropy does not behave like the discrete kind. It can be negative --- a
narrow Gaussian has negative differential entropy --- and it changes under a
change of units, since rescaling $x$ shifts $h$ by a constant. It is not ``the
amount of information'' in any absolute sense. Differences of entropies, which is
what every divergence below uses, are better behaved.
}

\whyhere{
Entropy is the building block for cross-entropy and KL divergence, which is how
the original GAN objective is expressed. It also gives the right language for mode
collapse: a generator producing one output has near-zero entropy, while the data
it was meant to imitate has high entropy. That gap is what the diversity guards in
this project detect.
}

%% ─────────────────────────────────────────────────────────────────────────────
\section{Cross-Entropy}
\label{sec:crossentropy}

\situation{
You design a code for weather reports, assuming rain and sun are equally likely,
so each gets a one-bit symbol. The real climate is rainy $90\%$ of the time. Your
code works, but wastes space --- you gave a common event as many bits as a rare
one.
}

\problem{
We need to measure the cost of using the \emph{wrong} distribution when the truth
is something else. That single quantity turns out to be the standard loss function
for classification, and it is what the original GAN discriminator minimises.
}

\workedarith{
Truth $p = (0.9, 0.1)$; assumed model $q = (0.5, 0.5)$.

\medskip
\textbf{Cost under the correct code} --- the entropy:
\[
  H(p) = 0.469 \text{ bits (from \S\ref{sec:entropy})} .
\]

\textbf{Cost under the wrong code} --- the cross-entropy:
\[
  H(p,q) = -\bigl[0.9\log_2 0.5 + 0.1\log_2 0.5\bigr] = -\log_2 0.5 = 1 \text{ bit} .
\]

\medskip
\textbf{The waste:} $1 - 0.469 = 0.531$ bits per report. That excess is the KL
divergence of \S\ref{sec:kl}.

\medskip
\textbf{A better model.} With $q = (0.8, 0.2)$:
\[
  H(p,q) = -\bigl[0.9\log_2 0.8 + 0.1\log_2 0.2\bigr]
         = -[0.9(-0.322) + 0.1(-2.322)] = 0.290 + 0.232 = 0.522 .
\]
Waste now $0.522 - 0.469 = 0.053$ bits. Closer model, smaller excess. \checkmark

\medskip
\textbf{As a classifier loss.} With a true label $y \in \{0,1\}$ and predicted
probability $\hat{y}$, binary cross-entropy is
\[
  \ell = -\bigl[y\ln\hat{y} + (1-y)\ln(1-\hat{y})\bigr] .
\]
At $y = 1$ and $\hat{y} = 0.9$: $\ell = -\ln 0.9 = 0.105$. At $\hat{y} = 0.5$:
$\ell = 0.693$. At $\hat{y} = 0.01$: $\ell = 4.605$. Confident and wrong is
punished heavily; the loss diverges as $\hat{y} \to 0$.
}

\formaldef{
The \textbf{cross-entropy} of $q$ relative to $p$ is
\[
  H(p,q) = -\sum_x p(x)\log q(x) ,
\]
the expected coding cost of using $q$ when the truth is $p$. It satisfies
$H(p,q) \geq H(p)$, with equality only when $q = p$ (Gibbs' inequality).

\textbf{Binary cross-entropy} is the two-outcome case and is the standard
classification loss.
}

\analogybreak{
Cross-entropy is not symmetric: $H(p,q) \neq H(q,p)$. It is a measure of the cost
of using $q$ to describe $p$, and swapping the roles is a different question. It
is also unbounded --- a single confident, wrong prediction contributes an
arbitrarily large loss, which makes it sensitive to mislabelled data.
}

\whyhere{
Binary cross-entropy is the discriminator's loss in the original GAN formulation
of Chapter~\ref{ch:gans}. Its unboundedness is part of the instability story:
when the discriminator becomes confident, the generator's gradient can vanish or
explode, which is one of the motivations for the Wasserstein alternative.
}

%% ─────────────────────────────────────────────────────────────────────────────
\section{KL Divergence}
\label{sec:kl}

\situation{
The wasted bits from the previous section --- the excess cost of the wrong code ---
deserve a name of their own, because that excess is exactly a measure of how far
apart two distributions are.
}

\problem{
We want a number that is zero when two distributions agree and grows as they
diverge. The coding waste supplies one, and it is the most widely used such
quantity --- but it has a defect that turns out to matter enormously here.
}

\workedarith{
\textbf{From the previous section.} $p = (0.9, 0.1)$, $q = (0.5, 0.5)$:
\[
  D_{\mathrm{KL}}(p\parallel q) = H(p,q) - H(p) = 1 - 0.469 = 0.531 \text{ bits} .
\]

\textbf{Directly:}
\[
  D_{\mathrm{KL}}(p\parallel q) = 0.9\log_2\frac{0.9}{0.5} + 0.1\log_2\frac{0.1}{0.5}
  = 0.9(0.848) + 0.1(-2.322) = 0.763 - 0.232 = 0.531 . \quad\checkmark
\]

\medskip
\textbf{Asymmetry, checked.}
\[
  D_{\mathrm{KL}}(q\parallel p) = 0.5\log_2\frac{0.5}{0.9} + 0.5\log_2\frac{0.5}{0.1}
  = 0.5(-0.848) + 0.5(2.322) = -0.424 + 1.161 = 0.737 .
\]
$0.531 \neq 0.737$. The two directions give different answers, so KL is not a
distance.

\medskip
\textbf{The fatal case.} Let $q$ assign probability zero where $p$ does not:
$p = (0.5, 0.5)$, $q = (1, 0)$. Then
\[
  D_{\mathrm{KL}}(p\parallel q) = 0.5\log_2\frac{0.5}{1} + 0.5\log_2\frac{0.5}{0} = \infty .
\]
Any region the model rules out but the data occupies makes KL infinite --- not
large, infinite. And an infinite loss has no usable gradient.
}

\formaldef{
The \textbf{Kullback--Leibler divergence} is
\[
  D_{\mathrm{KL}}(p\parallel q) = \sum_x p(x)\log\frac{p(x)}{q(x)}
  \qquad\text{or}\qquad
  \int p(x)\log\frac{p(x)}{q(x)}\,dx .
\]
Properties: $D_{\mathrm{KL}} \geq 0$, with equality only when $p = q$;
$D_{\mathrm{KL}}(p\parallel q) \neq D_{\mathrm{KL}}(q\parallel p)$; it violates the
triangle inequality; and it is $+\infty$ whenever $q(x) = 0$ at some $x$ with
$p(x) > 0$.
}

\analogybreak{
The two directions encode opposite preferences. Minimising
$D_{\mathrm{KL}}(p_{\text{data}}\parallel p_\theta)$ punishes the model for
ignoring any region the data occupies, producing over-spread,
everything-is-possible models. Minimising the reverse punishes putting mass where
the data has none, producing models that lock onto one mode and ignore the rest.
Mode-seeking versus mode-covering is a choice of direction, not a tuning
parameter.
}

\whyhere{
The infinity is the crux. A generator's output lies on a low-dimensional surface
in $\R^{128}$ (\S\ref{sec:datadist}), so early in training the model and the data
occupy essentially disjoint regions --- and KL is infinite everywhere. It reports
that the model is wrong without indicating any direction of improvement, which is
precisely what a gradient must supply.
}

%% ─────────────────────────────────────────────────────────────────────────────
\section{Jensen--Shannon Divergence and Disjoint Support}
\label{sec:jsd}

\situation{
Two islands with no bridge. Asked how to get from one to the other, an unhelpful
guide says only ``you cannot''. True, and useless --- you wanted a direction to
start walking, not a verdict.
}

\problem{
JS divergence fixes KL's asymmetry and its infinities, and is what the original
GAN objective implicitly minimises. It nonetheless fails in the same situation, for
a related reason: it becomes \emph{constant}, and a constant has zero gradient.
}

\workedarith{
\textbf{Two point masses that do not overlap.} Let $p$ put all its mass at $0$ and
$q$ all its mass at $\theta \neq 0$. Their mixture $m = \tfrac12(p+q)$ puts $1/2$
at each location.
\[
  D_{\mathrm{KL}}(p\parallel m) = 1 \cdot \log_2\frac{1}{1/2} = 1 ,
\]
and likewise $D_{\mathrm{KL}}(q\parallel m) = 1$. So
\[
  \mathrm{JSD}(p,q) = \tfrac12(1) + \tfrac12(1) = 1 \text{ bit} .
\]

\medskip
\textbf{Now move them closer.} Set $\theta = 100$, then $\theta = 10$, then
$\theta = 0.001$. As long as the supports do not overlap, the calculation is
identical:

\medskip
\begin{tabular}{@{}lrr@{}}
\toprule
$\theta$ & $\mathrm{JSD}$ & $W_1$ \\
\midrule
$100$   & $1$ & $100$ \\
$10$    & $1$ & $10$ \\
$1$     & $1$ & $1$ \\
$0.001$ & $1$ & $0.001$ \\
$0$     & $0$ & $0$ \\
\bottomrule
\end{tabular}

\medskip
\textbf{Read the JSD column.} It is $1$ everywhere and then drops discontinuously
to $0$. Its derivative with respect to $\theta$ is zero everywhere it is defined.
Gradient descent on JSD gets no signal at all --- the objective is flat, so the
generator has no indication that $\theta = 0.001$ is nearly right and
$\theta = 100$ is badly wrong.

\medskip
\textbf{Read the $W_1$ column.} It equals $|\theta|$: smooth, non-zero derivative
everywhere, pointing directly at the answer. This single comparison is the entire
argument for Wasserstein GANs.
}

\formaldef{
The \textbf{Jensen--Shannon divergence} is
\[
  \mathrm{JSD}(p,q) = \tfrac12 D_{\mathrm{KL}}(p\parallel m) + \tfrac12 D_{\mathrm{KL}}(q\parallel m) ,
  \qquad m = \tfrac12(p+q) .
\]
It is symmetric, always finite, bounded above by $1$ bit, and $\sqrt{\mathrm{JSD}}$
is a genuine metric.

Its defect: when $p$ and $q$ have disjoint support, $\mathrm{JSD} = \log 2$
regardless of how far apart they are, so its gradient with respect to the
generator's parameters is zero.
}

\analogybreak{
It is tempting to dismiss disjoint support as a pathological edge case. It is the
normal situation. A generator mapping $\R^{100}$ into $\R^{128}$ produces output
confined to a set of measure zero, so the model and data distributions are almost
surely disjoint at initialisation and remain so until they nearly coincide.
Adding noise to both restores overlap and is a real remedy, at the cost of blurring
what is being matched.
}

\whyhere{
This is the theoretical case for WGAN, and it is why QuantGAN and FinGAN in this
project use a WGAN-GP objective rather than the original cross-entropy formulation.
Chapter~\ref{ch:gans} sets out the practical consequences --- and the ways the
Wasserstein formulation fails to deliver everything the theory promises.
}

%% ─────────────────────────────────────────────────────────────────────────────
\section{Optimal Transport and the Wasserstein Distance}
\label{sec:wasserstein}

\situation{
Two piles of sand on a table, one the distribution of real returns and the other
of synthetic returns. How different are they? Ask for the minimum work --- mass
times distance --- needed to reshape the first pile into the second.
}

\problem{
KL and JS compare distributions \emph{pointwise}, asking how the densities differ
at each location, which is why they break when the supports do not meet.
Transporting mass instead uses the geometry of the underlying space, so
distributions that are close in position are close in value even when they never
overlap.
}

\workedarith{
\textbf{One dimension makes it easy.} Sort both samples and pair by rank; $W_1$ is
the average absolute gap.

Real $\{1, 3, 5\}$, synthetic $\{2, 3, 4\}$:
\[
  W_1 = \tfrac13\bigl(|1-2| + |3-3| + |5-4|\bigr) = \tfrac13(1+0+1) = 0.667 .
\]

\medskip
\textbf{Why sorted pairing is optimal.} Any crossing pair can be uncrossed without
increasing the cost. Pairing $1\!\to\!4$ and $5\!\to\!2$ costs $3 + 3 = 6$;
uncrossing to $1\!\to\!2$ and $5\!\to\!4$ costs $1 + 1 = 2$. Uncrossing never
hurts, so the sorted matching is optimal. \checkmark

\medskip
\textbf{The two point masses again.} $p$ at $0$, $q$ at $\theta$: all the mass
moves a distance $\theta$, so $W_1 = |\theta|$ --- matching the table in
\S\ref{sec:jsd} and confirming the gradient is $\pm 1$, never zero.

\medskip
\textbf{The shuffled control.} $W_1 = 0.000\text{e}{+}00$ exactly. A permutation
of the real data sorts to the identical sequence, so every paired gap is zero.

\medskip
\textbf{A floor that matters.} Suppose a generator collapses to emitting the
constant $\mu$. Then $W_1$ becomes the mean absolute deviation of the real
data --- for BOVESPA, $0.00770$. That is the \emph{worst} score a
constant-output model can achieve on this metric, and a genuine but imperfect
generator can easily do worse than $0.00770$. On $W_1$ alone, a collapsed model
can beat a working one.
}

\formaldef{
The \textbf{Wasserstein-$p$ distance} is
\[
  W_p(P,Q) = \left(\inf_{\gamma\in\Gamma(P,Q)} \E_{(X,Y)\sim\gamma}\norm{X-Y}^{p}\right)^{1/p} ,
\]
the infimum over \textbf{couplings} $\gamma$ --- joint distributions whose
marginals are $P$ and $Q$. Each coupling is a transport plan; $\gamma(x,y)$ is the
mass moved from $x$ to $y$.

In one dimension with $p=1$ this collapses to
\[
  W_1(P,Q) = \int_{-\infty}^{\infty}\bigl|F_P(x) - F_Q(x)\bigr|\,dx
           = \int_0^1\bigl|F_P^{-1}(u) - F_Q^{-1}(u)\bigr|\,du .
\]
$W_p$ is a genuine metric: symmetric, zero only when $P = Q$, and satisfying the
triangle inequality.
}

\analogybreak{
$W_1$ is permutation-invariant, so $W_1(\text{real},\text{shuffled}) = 0$. That is
correct --- the two have identical distributions --- and a hard limitation, since
it is blind to ordering by construction. No refinement of the computation will make
it detect a shuffle, which is why it belongs in the fidelity family and why the
evaluation needs a second family alongside it.

The optimisation is also genuinely hard above one dimension: the infimum over
couplings has no closed form and exact solution costs roughly $O(n^3)$. One
dimension is the lucky case, and it is the case this project's metric occupies.
}

\whyhere{
\texttt{wasserstein} is one of the seven ranked fidelity metrics, computed on the
sorted one-dimensional samples. The MAD floor above is the reason it cannot be
read in isolation: Chapter~\ref{ch:evaluation} shows how the composite design and
the degeneracy guards together prevent a collapsed model from winning on it.
}

%% ─────────────────────────────────────────────────────────────────────────────
\section{Kantorovich--Rubinstein Duality}
\label{sec:duality}

\situation{
Two ways to price a removals job. Either plan every item's journey and total the
work --- hard, with a vast number of possible plans. Or ask what the most
aggressive honest valuer would charge, given that their pricing cannot be
arbitrarily steep. Remarkably, the two answers agree.
}

\problem{
The definition of $W_1$ is an infimum over all couplings, which is not something a
neural network can compute. To train against Wasserstein distance we need a form
involving a maximisation over \emph{functions} instead --- because maximising over
functions is exactly what a network does.
}

\workedarith{
\textbf{The dual form.}
\[
  W_1(P,Q) = \sup_{\norm{f}_L \leq 1}\Bigl(\E_{x\sim P}[f(x)] - \E_{y\sim Q}[f(y)]\Bigr) ,
\]
where the supremum runs over all $1$-Lipschitz functions: those with
$|f(a)-f(b)| \leq |a-b|$ for all $a,b$.

\medskip
\textbf{Check it on the two point masses.} $P$ at $0$, $Q$ at $\theta > 0$. Then
\[
  \E_P[f] - \E_Q[f] = f(0) - f(\theta) .
\]
The Lipschitz condition caps this: $|f(0)-f(\theta)| \leq |\theta|$. The bound is
attained by $f(x) = -x$, which is exactly $1$-Lipschitz and gives
$f(0)-f(\theta) = \theta$. So the supremum is $\theta = W_1$. \checkmark

\medskip
\textbf{Why the constraint is essential.} Without it, take $f(x) = -Kx$ for large
$K$: the difference becomes $K\theta$, unbounded. The supremum would be $+\infty$
for any two distinct distributions and would measure nothing. The Lipschitz bound
is what makes the dual finite and meaningful.

\medskip
\textbf{A three-point check.} $P$ uniform on $\{0,2\}$, $Q$ uniform on $\{1,3\}$.
Primal: pair $0\!\to\!1$ and $2\!\to\!3$, each distance 1, so $W_1 = 1$.
Dual with $f(x) = -x$:
\[
  \E_P[f] - \E_Q[f] = \tfrac12(0 + (-2)) - \tfrac12((-1) + (-3)) = -1 - (-2) = 1 . \quad\checkmark
\]
}

\formaldef{
A function $f$ is \textbf{$K$-Lipschitz} if
\[
  |f(a) - f(b)| \leq K\,\norm{a-b} \quad\text{for all } a, b .
\]
For differentiable $f$ this is equivalent to $\norm{\nabla f} \leq K$ everywhere.

\textbf{Kantorovich--Rubinstein duality:}
\[
  W_1(P,Q) = \sup_{\norm{f}_L\leq 1}\Bigl(\E_{P}[f] - \E_{Q}[f]\Bigr) .
\]
The maximising $f$ is called the \textbf{Kantorovich potential}; in a GAN it is
approximated by the \textbf{critic}.
}

\analogybreak{
Duality is exact; the neural approximation is not. A network with a Lipschitz
constraint is not the full space of $1$-Lipschitz functions, so the supremum is
taken over a restricted family and the result is a lower bound on $W_1$. Nor is
the supremum reached: the critic is trained for a finite number of steps and
merely approaches it. What the generator descends is therefore an estimate of an
estimate --- a gap Chapter~\ref{ch:gans} takes seriously.
}

\whyhere{
This is where $\norm{\nabla D}_2 = 1$ comes from, and every piece is now in place.
Enforcing $1$-Lipschitz continuity means bounding the critic's gradient norm by 1
(\S\ref{sec:norms}). WGAN-GP enforces it softly, as the penalty of
\S\ref{sec:lagrange}, adding
\[
  \lambda\,\E\Bigl[\bigl(\norm{\nabla D(\hat{x})}_2 - 1\bigr)^2\Bigr]
\]
with $\lambda = 10$ throughout this project. The target is $1$ rather than
``at most 1'' because the optimal critic saturates the constraint --- as
$f(x) = -x$ did above, achieving exactly slope 1 where it matters.
}

%% ─────────────────────────────────────────────────────────────────────────────
\section{Diffusion Models and Score Matching}
\label{sec:diffusion}

\situation{
A drop of ink in still water spreads until the glass is uniformly grey. Running
that backwards --- grey water reassembling into a drop --- is impossible in
physics. But if you had recorded, at every instant, which way the ink was
spreading, you could retrace it.
}

\problem{
GANs avoid likelihoods at the cost of adversarial instability
(\S\ref{sec:minimax}). An alternative keeps a tractable training objective and no
second player at all, by learning to reverse a noising process step by step.
}

\workedarith{
\textbf{Forward: destroy the data.} Add Gaussian noise in small steps:
\[
  x_t = \sqrt{1-\beta_t}\,x_{t-1} + \sqrt{\beta_t}\,\varepsilon_t ,
  \qquad \varepsilon_t \sim \N(0,I) .
\]
With $\beta_t = 0.01$ held fixed, the signal retained after $t$ steps is
$(1-0.01)^{t/2}$:

\medskip
\begin{tabular}{@{}lrr@{}}
\toprule
$t$ & signal retained $0.99^{t/2}$ & \\
\midrule
$10$    & $0.951$ & mostly intact \\
$100$   & $0.605$ & degraded \\
$1000$  & $0.0067$ & essentially pure noise \\
\bottomrule
\end{tabular}

\medskip
After $1{,}000$ steps the data is gone and the distribution is $\N(0,I)$ --- which
we can sample from trivially.

\medskip
\textbf{Backward: learn to undo one step.} Train a network $\varepsilon_\theta$ to
predict the noise that was added:
\[
  L = \E_{t,x_0,\varepsilon}\Bigl[\norm{\varepsilon - \varepsilon_\theta(x_t, t)}^2\Bigr] .
\]
This is an ordinary squared-error regression. No adversary, no minimax, no
Lipschitz constraint --- just a target the network can be scored against directly.

\medskip
\textbf{Generation.} Start from pure noise and apply the learned reverse step
$1{,}000$ times. The cost is $1{,}000$ network evaluations per sample against a
GAN's one, which is the central practical trade.
}

\formaldef{
A \textbf{denoising diffusion probabilistic model} (Ho et al.\ 2020) defines a
forward chain adding Gaussian noise over $T$ steps and learns the reverse. The
training objective is a variational bound on the log-likelihood that reduces to
the simple noise-prediction regression above.

The \textbf{score} of a density is $\nabla_x \log p(x)$, and
\textbf{score matching} learns it directly; predicting the added noise is
equivalent to estimating the score up to a scale factor. Sampling then follows the
score uphill toward regions of high density.
}

\analogybreak{
Diffusion models trade instability for cost. Training is stable --- an ordinary
regression --- and sampling is one to three orders of magnitude slower than a GAN,
since it requires hundreds or thousands of forward passes. They also assume the
Gaussian noising process is appropriate to the data, which for heavy-tailed
financial returns is a real question rather than a formality: the forward process
attenuates precisely the extreme values that matter most here.
}

\whyhere{
Diffusion-TS is the intended fourth baseline for this benchmark and is not yet
implemented. The evaluation framework is deliberately model-agnostic --- it scores
samples, never likelihoods or internals --- so adding a diffusion model requires no
change to the metrics. The parity question of \S\ref{sec:sgd} will need care,
though: ``equal generator updates'' is well defined for the three GANs, and a
diffusion model has no adversarial generator step, so a defensible common budget
would have to be argued rather than assumed.
}
